\section{模式五：环境隔离与可复现性}

\subsection{环境对 Deep Agent 评测的关键影响}

Deep Agent 具有状态性，能够读写文件、调用 API、修改系统配置。这意味着如果测试环境没有正确隔离，\textbf{前一个测试用例的副作用会污染后续测试}，导致难以复现的测试失败（flaky tests）。

\begin{importantbox}{Deep Agent 评测的黄金法则}
每个 eval 运行必须获得一个\textbf{全新、干净的环境}，以确保结果可复现。这是 Deep Agent 评测与传统 LLM 评测最大的基础设施差异。
\end{importantbox}

\subsection{编程 Agent 的环境隔离策略}

文章以编程 Agent 为例介绍了两种隔离策略：

\begin{table}[H]
\centering
\caption{环境隔离策略对比}
\begin{tabular}{lp{4.5cm}p{4.5cm}}
\toprule
\textbf{特性} & \textbf{Docker 容器/Sandbox} & \textbf{临时目录} \\
\midrule
隔离程度 & 完全隔离（进程、文件系统、网络） & 文件系统级隔离 \\
启动开销 & 较高（秒级） & 极低（毫秒级） \\
适用场景 & 需要网络隔离或系统级操作的 Agent & 仅需文件系统操作的 Agent \\
代表实现 & Harbor / TerminalBench & DeepAgents CLI 临时目录方案 \\
\bottomrule
\end{tabular}
\end{table}

\begin{knowledgebox}{Harbor 与 TerminalBench}
Harbor 是一个为编程 Agent 评测设计的平台，TerminalBench 是其提供的评测基准之一。它为每个评测任务启动独立的 Docker 容器，确保 Agent 的文件操作、进程管理和网络访问互不干扰。对于生产级 Agent 评测系统，这类重量级隔离方案是推荐做法。
\end{knowledgebox}

\subsection{Mock API 请求：更快、更可控}

除了文件系统隔离，API 调用也需要考虑。LangSmith Assist 需要连接真实的 LangSmith API，但直接调用生产 API 存在两个问题：

\begin{enumerate}
  \item \textbf{速度慢}：网络请求增加延迟，拖慢整个评测流程
  \item \textbf{成本高}：大量 API 调用可能产生费用
  \item \textbf{不稳定}：外部服务的状态变化会影响测试结果的可复现性
\end{enumerate}

解决方案是\textbf{录制-回放（record-replay）模式}：

\begin{table}[H]
\centering
\caption{API Mock 工具推荐}
\begin{tabular}{llp{7cm}}
\toprule
\textbf{语言} & \textbf{工具} & \textbf{原理} \\
\midrule
Python & VCR.py & 首次运行录制 HTTP 交互到文件，后续回放 \\
JavaScript & Hono proxy & 通过代理应用拦截并回放 fetch 请求 \\
\bottomrule
\end{tabular}
\end{table}

\begin{warningbox}{Mock 的隐患：API 变更导致过期}
录制-回放模式的一个风险是当上游 API 发生变更时，录制的响应可能已经过时（stale）。建议定期用真实 API 重新录制一次基准数据，并在 CI 中设置周期性的``真实 API''测试作为兜底。
\end{warningbox}

\subsection{本章小结}

环境隔离和可复现性是 Deep Agent 评测的基础设施层保障。文件系统操作需要临时目录或容器隔离，API 调用需要录制-回放机制。两者结合可以显著提高评测的速度、成本效率和结果稳定性。


